\documentclass[10pt,a4paper]{amsart}
\usepackage{amsmath,amssymb,amsthm,graphicx}
% Retained preamble for the Stacks Project chapter morphisms.tex (Morphisms of Schemes).
% The chapter file pinned-original-source.tex carries only \input{preamble} on its first line, so no
% preamble text lives inside the chapter itself.  The Stacks Project uses one single shared preamble for
% every chapter; the snapshot stages it verbatim in sources/stacks-chow/preamble.tex.  The block below is
% copied mechanically, line for line, from that file's lines 14-18 (the Xy-pic setup) and lines 154-217
% (the native theorem declarations together with the mathematical macros).  Nothing was added, removed or
% redefined.  The wrapper class and the amsmath/amssymb/amsthm/graphicx base packages are supplied by the
% assembler, so the class/\externaldocument part of the shared preamble is not repeated here.
% For commutative diagrams we use Xy-pic
\usepackage[all]{xy}

% We use 2cell for 2-commutative diagrams.
\xyoption{2cell}
% Theorem environments.
%
\theoremstyle{plain}
\newtheorem{theorem}[subsection]{Theorem}
\newtheorem{proposition}[subsection]{Proposition}
\newtheorem{lemma}[subsection]{Lemma}

\theoremstyle{definition}
\newtheorem{definition}[subsection]{Definition}
\newtheorem{example}[subsection]{Example}
\newtheorem{exercise}[subsection]{Exercise}
\newtheorem{situation}[subsection]{Situation}

\theoremstyle{remark}
\newtheorem{remark}[subsection]{Remark}
\newtheorem{remarks}[subsection]{Remarks}

\numberwithin{equation}{subsection}

% Macros
%
\def\lim{\mathop{\mathrm{lim}}\nolimits}
\def\colim{\mathop{\mathrm{colim}}\nolimits}
\def\Spec{\mathop{\mathrm{Spec}}}
\def\Hom{\mathop{\mathrm{Hom}}\nolimits}
\def\Ext{\mathop{\mathrm{Ext}}\nolimits}
\def\SheafHom{\mathop{\mathcal{H}\!\mathit{om}}\nolimits}
\def\SheafExt{\mathop{\mathcal{E}\!\mathit{xt}}\nolimits}
\def\Sch{\mathit{Sch}}
\def\Mor{\mathop{\mathrm{Mor}}\nolimits}
\def\Ob{\mathop{\mathrm{Ob}}\nolimits}
\def\Sh{\mathop{\mathit{Sh}}\nolimits}
\def\NL{\mathop{N\!L}\nolimits}
\def\CH{\mathop{\mathrm{CH}}\nolimits}
\def\proetale{{pro\text{-}\acute{e}tale}}
\def\etale{{\acute{e}tale}}
\def\QCoh{\mathit{QCoh}}
\def\Ker{\mathop{\mathrm{Ker}}}
\def\Im{\mathop{\mathrm{Im}}}
\def\Coker{\mathop{\mathrm{Coker}}}
\def\Coim{\mathop{\mathrm{Coim}}}

% Boxtimes
%
\DeclareMathSymbol{\boxtimes}{\mathbin}{AMSa}{"02}

%
% Macros for moduli stacks/spaces
%
\def\QCohstack{\mathcal{QC}\!\mathit{oh}}
\def\Cohstack{\mathcal{C}\!\mathit{oh}}
\def\Spacesstack{\mathcal{S}\!\mathit{paces}}
\def\Quotfunctor{\mathrm{Quot}}
\def\Hilbfunctor{\mathrm{Hilb}}
\def\Curvesstack{\mathcal{C}\!\mathit{urves}}
\def\Polarizedstack{\mathcal{P}\!\mathit{olarized}}
\def\Complexesstack{\mathcal{C}\!\mathit{omplexes}}
% \Pic is the operator that assigns to X its picard group, usage \Pic(X)
% \Picardstack_{X/B} denotes the Picard stack of X over B
% \Picardfunctor_{X/B} denotes the Picard functor of X over B
\def\Pic{\mathop{\mathrm{Pic}}\nolimits}
\def\Picardstack{\mathcal{P}\!\mathit{ic}}
\def\Picardfunctor{\mathrm{Pic}}
\def\Deformationcategory{\mathcal{D}\!\mathit{ef}}

% --- appended 2026-10-12 (Chief): the project own definitions of the reference/slogan/history
% display environments, copied mechanically from the shared Stacks preamble lines 8-12
% (\usepackage{verbatim} plus the three \newenvironment lines, which the chapter preamble keeps
% verbatim from the project preamble). Nothing added, removed or redefined: these environments
% carry no mathematical content and the project itself typesets them as comments. Without these
% three lines a counted statement whose own text uses them fails with Environment slogan undefined.
\usepackage{verbatim}
\newenvironment{reference}{\comment}{\endcomment}
\newenvironment{slogan}{\comment}{\endcomment}
\newenvironment{history}{\comment}{\endcomment}
\usepackage{cleveref}
\crefname{originallemma}{original lemma}{original lemmas}
\Crefname{originallemma}{Original lemma}{Original lemmas}
\makeatletter
\newlabel{algebra-lemma-flat}{{00HD}{source}{}{sourcepointer-1}{}}
\newlabel{algebra-lemma-flat@cref}{{[originallemma][0][]00HD}{[1][0][]source}{}{}{}}
\makeatother
\crefname{originallemma}{original lemma}{original lemmas}
\Crefname{originallemma}{Original lemma}{Original lemmas}
\makeatletter
\newlabel{algebra-lemma-flatness-descends-more-general}{{0584}{source}{}{sourcepointer-2}{}}
\newlabel{algebra-lemma-flatness-descends-more-general@cref}{{[originallemma][0][]0584}{[1][0][]source}{}{}{}}
\makeatother
\crefname{originallemma}{original lemma}{original lemmas}
\Crefname{originallemma}{Original lemma}{Original lemmas}
\makeatletter
\newlabel{algebra-lemma-local-flat-ff}{{00HR}{source}{}{sourcepointer-3}{}}
\newlabel{algebra-lemma-local-flat-ff@cref}{{[originallemma][0][]00HR}{[1][0][]source}{}{}{}}
\makeatother
\crefname{originallemma}{original lemma}{original lemmas}
\Crefname{originallemma}{Original lemma}{Original lemmas}
\makeatletter
\newlabel{sheaves-lemma-stalk-pullback-modules}{{0098}{source}{}{sourcepointer-4}{}}
\newlabel{sheaves-lemma-stalk-pullback-modules@cref}{{[originallemma][0][]0098}{[1][0][]source}{}{}{}}
\makeatother
\title{Flatness over the base is detected after pulling back along a flat morphism}
\author{The Stacks Project Authors}
\date{stacks-project:morphisms.tex:local-snapshot-20261012; GNU FDL 1.2 or later}
\begin{document}
\maketitle
\noindent\textit{Editorial note.} All delivered mathematical text is the author's own, copied from the pinned Stacks Project chapter file pinned-original-source.tex; no statement or proof step was rewritten, reordered or removed.\par\medskip
\noindent\textit{Editorial note.} Following the multi-claim convention for this source, source\_url is the source-level URL https://stacks.math.columbia.edu/ and the source label lemma-flat-permanence is carried in claim\_key and here rather than in a per-claim tag URL.\par\medskip
\noindent\textit{Editorial note.} The preamble is the chapter's shared Stacks Project preamble, retained as sources/stacks-morphisms/retained-preamble.tex: a line-for-line excerpt of the shared Stacks preamble staged with this snapshot. No macro was invented or redefined, and this card adds no preamble text of its own.\par\medskip
\noindent\textit{Editorial note.} Statement span 4880-4894 and proof span 4896-4925 come from the halfA pool card operations/pool-cards/stacks-morphisms-halfA.json and were re-read in the pinned chapter file before this card was written.\par\medskip
\noindent\textit{Editorial note.} The proof opens with the author's remark that the statement can also be proved by a cited algebra lemma, followed by 'Here is a direct proof.'; both halves are preserved and the direct proof is the counted argument.\par\medskip
\noindent\textit{Editorial note.} 4 cross-references occur inside the counted spans; they are listed in external\_references as navigation pointers only.\par\medskip
\noindent\textit{Allowed standard prerequisites (editorial).}\begin{itemize}
\item A ring map that is flat at a prime is faithfully flat after localizing at that prime, cited from the algebra chapter. Cited at proof lines 5116-5117. Makes the local ring map from the base point to the source point faithfully flat, which is what allows the kernel computation to be tested after tensoring.
\item The stalk of a pullback module is the tensor product of the stalk with the local ring, cited from the sheaves chapter. Cited at proof lines 5119-5120. Identifies the stalk at the point of the pulled back sheaf with the tensor product used in the kernel computation.
\item The characterization of flatness by the vanishing of the kernel of the tensor product with an injection of modules, cited from the algebra chapter. Cited at proof lines 5140-5142. Converts the displayed equality of kernels into the equivalence of flatness statements asserted by the lemma.
\end{itemize}\medskip
\ifdefined\hypertarget\hypertarget{sourcepointer-1}{}\fi
\noindent\textit{Source reference.} Original lemma 00HD: focused-seeds/AGP/sources/stacks-snapshot-20261011/tags.txt, source lines 639--639. Same-chapter reference used in the retained author text at source line(s) 4924. It is supplied only as a navigation pointer so that the retained text compiles; the counted item delivers the author's own statement and the author's complete proof, and no part of the referenced item is copied in or substituted for the counted argument.\par\medskip
\ifdefined\hypertarget\hypertarget{sourcepointer-2}{}\fi
\noindent\textit{Source reference.} Original lemma 0584: focused-seeds/AGP/sources/stacks-snapshot-20261011/tags.txt, source lines 6767--6767. Same-chapter reference used in the retained author text at source line(s) 4898. It is supplied only as a navigation pointer so that the retained text compiles; the counted item delivers the author's own statement and the author's complete proof, and no part of the referenced item is copied in or substituted for the counted argument.\par\medskip
\ifdefined\hypertarget\hypertarget{sourcepointer-3}{}\fi
\noindent\textit{Source reference.} Original lemma 00HR: focused-seeds/AGP/sources/stacks-snapshot-20261011/tags.txt, source lines 653--653. Same-chapter reference used in the retained author text at source line(s) 4904. It is supplied only as a navigation pointer so that the retained text compiles; the counted item delivers the author's own statement and the author's complete proof, and no part of the referenced item is copied in or substituted for the counted argument.\par\medskip
\ifdefined\hypertarget\hypertarget{sourcepointer-4}{}\fi
\noindent\textit{Source reference.} Original lemma 0098: focused-seeds/AGP/sources/stacks-snapshot-20261011/tags.txt, source lines 344--344. Same-chapter reference used in the retained author text at source line(s) 720, 4907. It is supplied only as a navigation pointer so that the retained text compiles; the counted item delivers the author's own statement and the author's complete proof, and no part of the referenced item is copied in or substituted for the counted argument.\par\medskip
\begin{lemma}
\label{lemma-flat-permanence}
Let $h : X \to Y$ be a morphism of schemes over $S$.
Let $\mathcal{G}$ be a quasi-coherent sheaf on $Y$.
Let $x \in X$ with $y = h(x) \in Y$. If $h$ is flat at $x$, then
$$
\mathcal{G}\text{ flat over }S\text{ at }y
\Leftrightarrow
h^*\mathcal{G}\text{ flat over }S\text{ at }x.
$$
In particular: If $h$ is surjective and flat, then
$\mathcal{G}$ is flat over $S$, if and only if
$h^*\mathcal{G}$ is flat over $S$. If $h$ is surjective and
flat, and $X$ is flat over $S$, then $Y$ is flat over $S$.
\end{lemma}

\begin{proof}
You can prove this by applying
Algebra, Lemma \ref{algebra-lemma-flatness-descends-more-general}.
Here is a direct proof. Let $s \in S$ be the image of $y$.
Consider the local ring maps
$\mathcal{O}_{S, s} \to \mathcal{O}_{Y, y} \to \mathcal{O}_{X, x}$.
By assumption the ring map $\mathcal{O}_{Y, y} \to \mathcal{O}_{X, x}$
is faithfully flat, see
Algebra, Lemma \ref{algebra-lemma-local-flat-ff}.
Let $N = \mathcal{G}_y$. Note that
$h^*\mathcal{G}_x = N \otimes_{\mathcal{O}_{Y, y}} \mathcal{O}_{X, x}$, see
Sheaves, Lemma \ref{sheaves-lemma-stalk-pullback-modules}.
Let $M' \to M$ be an injection of $\mathcal{O}_{S, s}$-modules.
By the faithful flatness mentioned above we have
\begin{align*}
\Ker(
M' \otimes_{\mathcal{O}_{S, s}} N \to M \otimes_{\mathcal{O}_{S, s}} N)
\otimes_{\mathcal{O}_{Y, y}} \mathcal{O}_{X, x} \\
=
\Ker(
M' \otimes_{\mathcal{O}_{S, s}} N
\otimes_{\mathcal{O}_{Y, y}} \mathcal{O}_{X, x}
\to
M \otimes_{\mathcal{O}_{S, s}} N
\otimes_{\mathcal{O}_{Y, y}} \mathcal{O}_{X, x})
\end{align*}
Hence the equivalence of the lemma follows from the second characterization
of flatness in
Algebra, Lemma \ref{algebra-lemma-flat}.
\end{proof}

\end{document}
